\begin{abstract}
The design of timber bridges involves multiple design variables and conflicting code requirements, and is traditionally carried out through iterative procedures that depend on the designer's experience, without a systematic investigation of how span, strength class and dimensional variability interact in the structural behaviour.
This study proposes a parametric robust multi-objective optimization framework for roundwood timber bridges, with two conflicting objective functions, the timber volume and the utilization of the deflection limit, solved by NSGA-II under the ultimate and serviceability limit state constraints of ABNT NBR 7190:2022.
Robustness is incorporated by keeping the objectives at the nominal point and aggregating the constraints by the worst case over a deterministic five-point deviation grid $\pm\rho$ ($\rho = 5\%$) applied to the girder diameter.
The framework was applied to 20 configurations combining four spans (3.0 to 6.0~m) and the five strength classes D20 to D60 under the TB-240 design vehicle.
All configurations converged to feasible solutions, with timber consumption between 1.75 and 9.98~m\textsuperscript{3}.
Consumption grows with span according to a power law with exponent between 1.64 and 1.73, and the first class upgrade, from D20 to D30, accounts for 41 to 48\% of the total saving available up to D60.
Bending governs all configurations, fifteen through the deck and five through the girder, and none is governed by deflection, although deflection approaches its limit in the most severe cases.
Robustness carries a real and monotonic cost in volume, from $+6.0\%$ ($\rho=5\%$) to $+26.9\%$ ($\rho=20\%$).
The results were condensed into preliminary design charts and a pre-sizing equation ($R^2=0.999$), valid within the investigated domain, quantifying how span, strength class and dimensional variability control material demand and the governing limit states.
\end{abstract}


\textbf{Keywords}: roundwood timber bridges, multi-objective optimization, robust design, parametric design, structural optimization, timber structures.
